\documentclass[11pt,a4paper]{article}

% ── Packages ──────────────────────────────────────────────────────────
\usepackage[utf8]{inputenc}
\usepackage[T1]{fontenc}
\usepackage{lmodern}
\usepackage[margin=2.5cm]{geometry}
\usepackage{booktabs}
\usepackage{tabularx}
\usepackage{longtable}
\usepackage{array}
\usepackage{xcolor}
\usepackage{listings}
\usepackage{hyperref}
\usepackage{fancyhdr}
\usepackage{titlesec}
\usepackage{enumitem}
\usepackage{amssymb}
\usepackage{amsmath}
\usepackage{graphicx}
\usepackage{float}
\usepackage{caption}
\usepackage{bytefield}
\bytefieldsetup{endianness=big}
\usepackage{multicol}
\usepackage{tcolorbox}
\usepackage{tikz}

% ── Colours ───────────────────────────────────────────────────────────
\definecolor{ee8blue}{RGB}{0,90,160}
\definecolor{ee8gray}{RGB}{100,100,100}
\definecolor{codebg}{RGB}{248,248,248}
\definecolor{codeborder}{RGB}{200,200,200}
\definecolor{keywordcolor}{RGB}{0,100,170}
\definecolor{commentcolor}{RGB}{100,140,100}
\definecolor{stringcolor}{RGB}{170,60,0}
\definecolor{warnbg}{RGB}{255,250,230}
\definecolor{warnborder}{RGB}{220,180,50}
\definecolor{notebg}{RGB}{235,245,255}
\definecolor{noteborder}{RGB}{70,130,200}
\definecolor{greenbg}{RGB}{235,250,240}
\definecolor{greenborder}{RGB}{60,160,100}

% ── Listings style ────────────────────────────────────────────────────
\lstdefinelanguage{ee8asm}{
  morekeywords={ADD,SUB,AND,OR,XOR,SHL,SHR,MOV,LDI,LT,EQ,JMP,JZ,JNZ,HALT,IN,OUT},
  sensitive=true,
  morecomment=[l]{;},
  morestring=[b]",
}

\lstset{
  language=ee8asm,
  basicstyle=\ttfamily\small,
  keywordstyle=\color{keywordcolor}\bfseries,
  commentstyle=\color{commentcolor}\itshape,
  stringstyle=\color{stringcolor},
  backgroundcolor=\color{codebg},
  frame=single,
  rulecolor=\color{codeborder},
  framesep=5pt,
  xleftmargin=10pt,
  xrightmargin=10pt,
  numbers=none,
  breaklines=true,
  showstringspaces=false,
  tabsize=4,
  columns=fullflexible,
  keepspaces=true,
}

% ── tcolorbox styles ─────────────────────────────────────────────────
\newtcolorbox{designnote}{
  colback=notebg,
  colframe=noteborder,
  fonttitle=\bfseries,
  title=Design Note,
  boxrule=0.5pt,
  arc=2pt,
  left=6pt, right=6pt, top=4pt, bottom=4pt,
}

\newtcolorbox{warningbox}{
  colback=warnbg,
  colframe=warnborder,
  fonttitle=\bfseries,
  title=Important,
  boxrule=0.5pt,
  arc=2pt,
  left=6pt, right=6pt, top=4pt, bottom=4pt,
}

\newtcolorbox{taskbox}[1][]{
  colback=greenbg,
  colframe=greenborder,
  fonttitle=\bfseries,
  title=#1,
  boxrule=0.5pt,
  arc=2pt,
  left=6pt, right=6pt, top=4pt, bottom=4pt,
}

\newtcolorbox{hintbox}{
  colback=codebg,
  colframe=codeborder,
  fonttitle=\bfseries,
  title=Hint,
  boxrule=0.5pt,
  arc=2pt,
  left=6pt, right=6pt, top=4pt, bottom=4pt,
}

% ── Hyperref setup ────────────────────────────────────────────────────
\hypersetup{
  colorlinks=true,
  linkcolor=ee8blue,
  urlcolor=ee8blue,
  citecolor=ee8blue,
  pdftitle={EE8 CPU -- Lab: Instruction Decoder and Control Unit},
  pdfauthor={EG2300 Course Team},
}

% ── Header / footer ──────────────────────────────────────────────────
\pagestyle{fancy}
\fancyhf{}
\fancyhead[L]{\small\textcolor{ee8gray}{EE8 CPU --- Lab: Instruction Decoder \& Control Unit}}
\fancyhead[R]{\small\textcolor{ee8gray}{EG2300}}
\fancyfoot[C]{\small\textcolor{ee8gray}{\thepage}}
\renewcommand{\headrulewidth}{0.4pt}
\renewcommand{\footrulewidth}{0pt}

% ── Section formatting ────────────────────────────────────────────────
\titleformat{\section}
  {\Large\bfseries\color{ee8blue}}{\thesection}{1em}{}[\vspace{-0.5em}\rule{\textwidth}{0.4pt}]
\titleformat{\subsection}
  {\large\bfseries\color{ee8blue!80!black}}{\thesubsection}{1em}{}
\titleformat{\subsubsection}
  {\normalsize\bfseries}{\thesubsubsection}{1em}{}

% ── Custom commands ───────────────────────────────────────────────────
\newcommand{\reg}[1]{\texttt{R#1}}
\newcommand{\opcode}[1]{\texttt{#1}}
\newcommand{\signal}[1]{\textsc{#1}}
\newcommand{\instr}[1]{\texttt{\textbf{#1}}}
\newcommand{\flag}[1]{\texttt{#1}}

% ══════════════════════════════════════════════════════════════════════
\begin{document}

% ── Title ─────────────────────────────────────────────────────────────
\begin{center}
{\Huge\bfseries\color{ee8blue} EE8 CPU}\\[0.4cm]
{\LARGE Lab: Instruction Decoder \& Control Unit}\\[0.8cm]
{\large Engineering Design Module (EG2300)}\\[0.3cm]
\rule{0.6\textwidth}{1pt}\\[0.5cm]
\end{center}

% ── Objectives ────────────────────────────────────────────────────────
\section{Objectives}

By the end of this lab you will have built a \textbf{combinational logic circuit} that takes a 4-bit opcode from the instruction word and produces the control signals (bit flags) that will drive every other component in your CPU.

\medskip
\noindent You will:
\begin{enumerate}[leftmargin=2em]
  \item Understand why a CPU needs control signals and what each one does.
  \item Extract the opcode from the 16-bit instruction word.
  \item Design and implement a combinational decoder that sets bit flags based on the opcode.
  \item Test your decoder against every opcode in the EE8 instruction set.
\end{enumerate}

% ── Background ────────────────────────────────────────────────────────
\section{Background}

\subsection{Where are we?}

So far you have built:
\begin{itemize}[leftmargin=2em]
  \item \textbf{PIPO registers} --- 8-bit storage from D flip-flops.
  \item \textbf{ROM} --- 16 rows of 16-bit instruction words.
  \item \textbf{Program Counter} --- T flip-flop counter that selects the current ROM address, with reset (all bits to 0) and an unconditional jump override.
\end{itemize}

Right now your ROM outputs a 16-bit instruction word every clock cycle, but nothing \emph{interprets} it. The instruction word is just a bundle of 16 wires carrying ones and zeros. The CPU needs logic that looks at those bits and decides:

\begin{itemize}[leftmargin=2em]
  \item Should a register be written to?
  \item Should the ALU add, subtract, or do something else?
  \item Is this a jump instruction?
  \item Should the CPU halt?
\end{itemize}

This is the job of the \textbf{instruction decoder} and \textbf{control unit}.

\subsection{What is a control signal?}

A \textbf{control signal} (or \textbf{bit flag}) is a single wire that is either 0 or 1.
It acts as an on/off switch for a specific part of the CPU.

\medskip
For example, you already have one control signal: the \textbf{Jump} flag on your program counter.
When \flag{Jump}~=~1, the PC loads the jump target address instead of incrementing.
When \flag{Jump}~=~0, the PC increments normally (PC~=~PC~+~1).

\medskip
The control unit's job is to produce \emph{all} of these flags simultaneously, based purely on the opcode. It is \textbf{combinational logic} --- no flip-flops, no clock, no memory. The opcode goes in, the flags come out.

\begin{designnote}
In a real CPU (ARM, RISC-V, x86) the control unit works the same way --- it decodes the instruction and asserts control signals that steer the datapath.
The only difference is scale: those CPUs have hundreds of signals; the EE8 has a handful.
You are building the same fundamental circuit.
\end{designnote}


% ── Part A: Field Extraction ──────────────────────────────────────────
\section{Part A: Extracting the Opcode}

The 16-bit instruction word from ROM has a specific structure.
Regardless of instruction type (R-type, I-type, or J-type), the \textbf{opcode is always bits [15:12]} --- the top 4 bits.

\begin{center}
\begin{bytefield}[bitwidth=1.3em]{16}
  \bitheader{0,1,2,3,4,5,6,7,8,9,10,11,12,13,14,15} \\
  \bitbox{4}{\small opcode} & \bitbox{12}{\small (varies by instruction type)}
\end{bytefield}
\end{center}

\begin{taskbox}[Task A]
Use a \textbf{splitter} in Logisim to extract bits [15:12] from the ROM output.
This gives you a 4-bit opcode signal.

\medskip
\noindent Your circuit should have:
\begin{itemize}[leftmargin=2em]
  \item \textbf{Input:} \texttt{instruction[15:0]} (16-bit wire from ROM)
  \item \textbf{Output:} \texttt{opcode[3:0]} (4-bit wire)
\end{itemize}
\end{taskbox}

You will extract the other fields (Rs1, Rs2, Rd, immediate, address) in a later lab when you build the register file and datapath. For now, the opcode is all the control unit needs.


% ── Part B: Core Control Signals ──────────────────────────────────────
\section{Part B: The Control Signals}

Your control unit will produce the following bit flags. Each flag is a single output wire (1 bit).

\subsection{Core Signals (Required)}

\renewcommand{\arraystretch}{1.4}
\begin{table}[H]
\centering
\begin{tabularx}{\textwidth}{l l X}
\toprule
\textbf{Signal} & \textbf{Width} & \textbf{Purpose} \\
\midrule
\flag{RegWrite} & 1 bit &
  \textbf{Register write enable.}
  When 1, the result of the current instruction will be written into a register.
  When 0, no register is modified.
  \emph{Example:} \instr{ADD} writes a result to Rd, so \flag{RegWrite}~=~1.
  \instr{JMP} does not write to any register, so \flag{RegWrite}~=~0. \\

\flag{ImmSel} & 1 bit &
  \textbf{Immediate select.}
  When 1, the data written to a register comes from the immediate field in the instruction (bits [7:0]) rather than from the ALU.
  Only the \instr{LDI} instruction uses this. \\

\flag{Jump} & 1 bit &
  \textbf{Unconditional jump.}
  When 1, the PC is loaded with the jump target address instead of PC+1.
  You already have this signal connected to your PC --- now you need to generate it from the opcode. \\

\flag{Branch} & 1 bit &
  \textbf{Conditional branch enable.}
  When 1, the PC \emph{may} be overridden with the branch target --- but only if the branch condition is met (checked against a register value).
  Used by \instr{JZ} and \instr{JNZ}. \\

\flag{Halt} & 1 bit &
  \textbf{Halt.}
  When 1, the PC stops incrementing and the CPU freezes.
  Only the \instr{HALT} instruction sets this. \\

\bottomrule
\end{tabularx}
\caption{Core control signals --- you must implement all of these.}
\end{table}

\begin{warningbox}
The \flag{Jump} signal you are generating here replaces the manual jump flag you may have used for testing previously.
After this lab, the jump flag will be \emph{automatically} set whenever the opcode is \opcode{1011} (\instr{JMP}).
\end{warningbox}

\subsection{Extension Signals (Optional)}

Once the core signals are working, you may add these for a more complete control unit.
These will become useful in later labs.

\renewcommand{\arraystretch}{1.4}
\begin{table}[H]
\centering
\begin{tabularx}{\textwidth}{l l X}
\toprule
\textbf{Signal} & \textbf{Width} & \textbf{Purpose} \\
\midrule
\flag{BranchCond} & 1 bit &
  Selects which branch condition to test.
  0~=~\instr{JZ} (jump if register is zero),
  1~=~\instr{JNZ} (jump if register is not zero).
  Only meaningful when \flag{Branch}~=~1. \\

\flag{ALUop[3:0]} & 4 bits &
  Tells the ALU which operation to perform (add, subtract, AND, etc.).
  For most R-type instructions this can be a direct passthrough of the opcode.
  You will connect this when you build the ALU. \\

\flag{FlagWrite} & 1 bit &
  When 1, the comparison result is written to \reg{3}.
  Only \instr{LT} and \instr{EQ} use this.
  Needed so that normal ALU instructions (\instr{ADD}, etc.) do not accidentally overwrite \reg{3}. \\

\flag{SysErr} & 1 bit &
  \textbf{Illegal instruction.}
  Set to 1 when the opcode is \opcode{1110} (reserved).
  Latched externally to halt the CPU and raise the alarm. \\

\bottomrule
\end{tabularx}
\caption{Extension control signals --- implement these for a more complete design.}
\end{table}


% ── Part C: Truth Table ───────────────────────────────────────────────
\section{Part C: The Truth Table}

The control unit is a \textbf{combinational logic function}: the 4-bit opcode is the input, and the control signals are the outputs. Before you build gates, you must first work out the truth table.

\begin{taskbox}[Task C1 --- Fill in the Truth Table]
Complete the table below. For each opcode, decide whether each control signal should be 0 or 1.

\medskip
\noindent Ask yourself for each instruction:
\begin{itemize}[leftmargin=2em]
  \item Does this instruction write a result to a register? $\rightarrow$ \flag{RegWrite}
  \item Does this instruction use an immediate value instead of the ALU? $\rightarrow$ \flag{ImmSel}
  \item Does this instruction jump unconditionally? $\rightarrow$ \flag{Jump}
  \item Does this instruction branch conditionally? $\rightarrow$ \flag{Branch}
  \item Does this instruction stop the CPU? $\rightarrow$ \flag{Halt}
\end{itemize}
\end{taskbox}

\renewcommand{\arraystretch}{1.3}
\begin{table}[H]
\centering
\begin{tabular}{c c l | c c c c c}
\toprule
\textbf{Opcode} & \textbf{Binary} & \textbf{Mnemonic} &
\rotatebox{60}{\flag{RegWrite}} &
\rotatebox{60}{\flag{ImmSel}} &
\rotatebox{60}{\flag{Jump}} &
\rotatebox{60}{\flag{Branch}} &
\rotatebox{60}{\flag{Halt}} \\
\midrule
0  & 0000 & ADD  & & & & & \\
1  & 0001 & SUB  & & & & & \\
2  & 0010 & AND  & & & & & \\
3  & 0011 & OR   & & & & & \\
4  & 0100 & XOR  & & & & & \\
5  & 0101 & SHL  & & & & & \\
6  & 0110 & SHR  & & & & & \\
7  & 0111 & MOV  & & & & & \\
8  & 1000 & LT   & & & & & \\
9  & 1001 & EQ   & & & & & \\
10 & 1010 & LDI  & & & & & \\
11 & 1011 & JMP  & & & & & \\
12 & 1100 & JZ   & & & & & \\
13 & 1101 & JNZ  & & & & & \\
14 & 1110 & (reserved) & & & & & \\
15 & 1111 & HALT & & & & & \\
\bottomrule
\end{tabular}
\caption{Control signal truth table --- fill this in before building your circuit.}
\label{tab:truth-table}
\end{table}

\begin{hintbox}
Look at the patterns in your completed table.
You should notice that several signals are only set for one or two opcodes, while others are set for a whole range.
This tells you how complex the gate logic needs to be for each signal.

\medskip
For example, \flag{Halt} is 1 for only \emph{one} opcode. That means you need a circuit that outputs 1 \emph{only} when the opcode is exactly \opcode{1111}.
How many gates does that take?
\end{hintbox}


% ── Part D: Implementation ────────────────────────────────────────────
\section{Part D: Build the Control Unit}

Now turn your truth table into a Logisim circuit using combinational logic gates.

\subsection{Approach}

For each output signal, you have a column of 0s and 1s in your truth table.
You need to design a logic circuit that produces that exact column given the 4-bit opcode input.

\medskip
\noindent There are two standard approaches:

\begin{enumerate}[leftmargin=2em]
  \item \textbf{Sum of products (SOP):} For each row where the output is 1, write a minterm (an AND of the 4 opcode bits or their complements). OR all the minterms together.
  This is what you would get from a Karnaugh map.

  \item \textbf{Direct pattern matching:} Some signals are 1 for a single opcode value.
  For these, you just need to detect that specific 4-bit pattern using AND and NOT gates.
  For example, \flag{Halt} is 1 only when opcode~=~\opcode{1111}, so:
  \[
    \flag{Halt} = \texttt{opcode[3]} \;\mathbin{\text{AND}}\; \texttt{opcode[2]} \;\mathbin{\text{AND}}\; \texttt{opcode[1]} \;\mathbin{\text{AND}}\; \texttt{opcode[0]}
  \]
  That is a single 4-input AND gate.
\end{enumerate}

\begin{designnote}
You may simplify your logic using Karnaugh maps, Boolean algebra, or by spotting patterns in the truth table. Any correct combinational implementation is acceptable.

\medskip
\noindent Some questions to consider while simplifying:
\begin{itemize}[leftmargin=2em]
  \item Are there signals where you only need to check one or two opcode bits?
  \item Are there groups of opcodes that share the same output pattern?
  \item Would it be simpler to compute a signal as NOT of something else?
\end{itemize}
\end{designnote}

\begin{taskbox}[Task D --- Build the Circuit]
Create a new Logisim subcircuit called \texttt{control\_unit}.

\medskip
\noindent \textbf{Interface:}
\begin{itemize}[leftmargin=2em]
  \item \textbf{Input:} \texttt{opcode[3:0]} (4 bits)
  \item \textbf{Outputs:} \flag{RegWrite}, \flag{ImmSel}, \flag{Jump}, \flag{Branch}, \flag{Halt} (1 bit each)
\end{itemize}

\medskip
\noindent Build the combinational logic for each output using AND, OR, NOT (and NAND, NOR, XOR if you wish).
\textbf{Do not use a ROM component} --- the point of this exercise is to design the logic yourself.

\medskip
\noindent If you are also implementing the extension signals, add them as additional outputs:
\flag{BranchCond}, \flag{ALUop[3:0]}, \flag{FlagWrite}, \flag{SysErr}.
\end{taskbox}


% ── Part E: Testing ───────────────────────────────────────────────────
\section{Part E: Testing Your Decoder}

\begin{taskbox}[Task E --- Verify All 16 Opcodes]
Test your control unit by manually setting the 4-bit opcode input to each value from \opcode{0000} to \opcode{1111} and checking that the output flags match your truth table.

\medskip
\noindent Use \textbf{probes} or \textbf{LEDs} on each output to see the flag values.

\medskip
\noindent Record your results in the table below (or take a screenshot for each opcode).
\end{taskbox}

\renewcommand{\arraystretch}{1.3}
\begin{table}[H]
\centering
\begin{tabular}{c l | c c c c c | c}
\toprule
\textbf{Opcode} & \textbf{Mnemonic} &
\rotatebox{60}{\flag{RegWrite}} &
\rotatebox{60}{\flag{ImmSel}} &
\rotatebox{60}{\flag{Jump}} &
\rotatebox{60}{\flag{Branch}} &
\rotatebox{60}{\flag{Halt}} &
\textbf{Pass?} \\
\midrule
0000 & ADD  & & & & & & $\square$ \\
0001 & SUB  & & & & & & $\square$ \\
0010 & AND  & & & & & & $\square$ \\
0011 & OR   & & & & & & $\square$ \\
0100 & XOR  & & & & & & $\square$ \\
0101 & SHL  & & & & & & $\square$ \\
0110 & SHR  & & & & & & $\square$ \\
0111 & MOV  & & & & & & $\square$ \\
1000 & LT   & & & & & & $\square$ \\
1001 & EQ   & & & & & & $\square$ \\
1010 & LDI  & & & & & & $\square$ \\
1011 & JMP  & & & & & & $\square$ \\
1100 & JZ   & & & & & & $\square$ \\
1101 & JNZ  & & & & & & $\square$ \\
1110 & (reserved) & & & & & & $\square$ \\
1111 & HALT & & & & & & $\square$ \\
\bottomrule
\end{tabular}
\caption{Test results --- fill in the observed outputs and tick if they match your truth table.}
\end{table}

\begin{hintbox}
If a signal is wrong for a specific opcode, trace back through your gates.
Check:
\begin{itemize}[leftmargin=2em]
  \item Are your NOT gates on the correct opcode bits?
  \item Did you accidentally swap two opcode bit positions in a splitter?
  \item Is the gate type correct (AND vs OR)?
\end{itemize}
\end{hintbox}


% ── Part F: Connecting to your PC ─────────────────────────────────────
\section{Part F: Connecting the Jump Flag to Your PC}

You already have a working program counter with an unconditional jump input.
Now you can connect the \flag{Jump} output from your control unit to that input, so the CPU automatically jumps when it encounters a \instr{JMP} instruction.

\begin{taskbox}[Task F --- Integration Test]
\begin{enumerate}[leftmargin=2em]
  \item Connect the ROM output through your splitter to extract \texttt{opcode[3:0]}.
  \item Feed the opcode into your control unit.
  \item Connect the \flag{Jump} output to your PC's jump-enable input.
  \item Connect the \flag{Halt} output to your PC's halt/freeze input (if you have one), or use it to gate the clock.
  \item Load a simple test program into your ROM:
\end{enumerate}

\begin{lstlisting}
; Test program: JMP should loop forever between addresses 0 and 2
0: LDI R0 42       ; 1010 00 00 00101010  -> 0xA02A
1: LDI R1 10       ; 1010 00 01 00001010  -> 0xA10A
2: JMP 0           ; 1011 00 000000 0000  -> 0xB000
3: HALT             ; 1111 000000000000   -> 0xF000
\end{lstlisting}

\noindent \textbf{Expected behaviour:}
\begin{itemize}[leftmargin=2em]
  \item The PC should count 0, 1, 2, 0, 1, 2, \ldots\ (looping).
  \item Address 3 (\instr{HALT}) should never be reached.
  \item The \flag{Jump} LED should light up only when the PC is at address 2.
  \item The \flag{RegWrite} LED should light up at addresses 0 and 1 (the \instr{LDI} instructions).
\end{itemize}
\end{taskbox}


% ── Deliverables ──────────────────────────────────────────────────────
\section{Deliverables}

\begin{enumerate}[leftmargin=2em]
  \item \textbf{Completed truth table} (Table~\ref{tab:truth-table}).
  \item \textbf{Logisim subcircuit:} \texttt{control\_unit} with correct inputs/outputs.
  \item \textbf{Test results} showing all 16 opcodes produce the correct flags.
  \item \textbf{Integration test} demonstrating the \flag{Jump} and \flag{Halt} signals working with your existing PC and ROM.
\end{enumerate}

\medskip
\noindent \textbf{Extension deliverables} (optional):
\begin{itemize}[leftmargin=2em]
  \item \flag{BranchCond} signal implemented and tested.
  \item \flag{ALUop[3:0]} output (can be a passthrough of the opcode for R-type instructions, don't-care for non-ALU instructions).
  \item \flag{FlagWrite} signal for \instr{LT}/\instr{EQ}.
  \item \flag{SysErr} signal for opcode \opcode{1110}.
\end{itemize}


% ── Summary ───────────────────────────────────────────────────────────
\section{Summary}

You have built the ``brain'' of the CPU --- the part that interprets instructions and tells every other component what to do. Every signal in your control unit is a simple combinational function of the 4-bit opcode.

\medskip
\noindent In the next lab, you will build the \textbf{register file}: four registers with multiplexers for reading (Rs1, Rs2) and a decoder for write-enable (Rd). The \flag{RegWrite} signal you built today will drive that write-enable logic.

\begin{designnote}
The instruction decoder pattern you used here --- extract a field, match patterns, produce control signals --- is the same pattern used in every processor ever made, from the 1971 Intel 4004 to modern ARM and RISC-V cores. The logic gets bigger, but the idea is identical.
\end{designnote}

\end{document}
